% GENERATED FILE. Edit canonical lessons, then run npm run generate:books.
% canonical-source-hash: fnv1a64:74e5da526e53df9d
% canonical-lessons: GE-C468-wut, GE-C468-zorn, GE-C468-aergern, GE-C468-aufregen, GE-C468-toben

\chapter{Wut und Zorn}
\label{ch:ge-wut-und-zorn}
\hlchaptermodality{468}

\begin{chapteropening}
\textbf{By the end of this chapter:} \emph{I can talk about anger and losing my temper.}

Complete the last lesson of chapter 468: I can talk about anger and losing my temper.
\end{chapteropening}

\section[die Wut]{die Wut — rage, fury}

\label{lesson:GE-C468-wut}



\begin{quote}
Before the new one: say the German for embarrassment, then the German for to blush.
\end{quote}

\subsection*{You'll want to know: die Wut}
\textbf{die Wut} — \textquotedblleft{}rage, fury\textquotedblright{}.

\begin{quote}
Vor Wut warf er sein Handy auf das Sofa.
\end{quote}

\begin{quote}
\emph{In a rage, he threw his phone onto the sofa.}
\end{quote}

No plural. \emph{Eine Wut haben auf} + accusative — to be furious with.

\begin{cousinweb}
\emph{Wut} is inherited from Middle High German \emph{wuot} and Old High German \emph{wuot} \textquotedblleft{}mad, madness\textquotedblright{}, from Proto-Germanic \emph{wōdaz} (EtymDB).
\end{cousinweb}

\begin{grammarlens}[title={What you've built}]
One more word for telling, explaining and describing at B1.
\end{grammarlens}

\subsection*{Your turn}
\begin{itemize}
\raggedright
  \item \emph{Say it:} \emph{die Wut}
  \item \emph{Say it:} \emph{die Wut}, once more
  \item \emph{Say it:} a sentence of your own with \emph{die Wut}
  \item \emph{Recall:} say the German for embarrassment, then the German for to blush, then say \emph{die Wut} again
\end{itemize}

\begin{tcolorbox}[breakable,colback=teal!4,colframe=teal!35!black,title={Before you move on}]
What does die Wut mean? (\textquotedblleft{}Rage, fury\textquotedblright{}.) Say it once more.
\end{tcolorbox}
\section[der Zorn]{der Zorn — anger, wrath}

\label{lesson:GE-C468-zorn}



\begin{quote}
Before the new one: say the German for to blush, then the German for rage.
\end{quote}

\subsection*{You'll want to know: der Zorn}
\textbf{der Zorn} — \textquotedblleft{}anger, wrath\textquotedblright{}.

\begin{quote}
Der Zorn des Vaters war schnell vorbei.
\end{quote}

\begin{quote}
\emph{The father's anger was soon over.}
\end{quote}

No plural. More literary than \emph{die Wut}.

\begin{cousinweb}
\emph{Zorn} is inherited from Middle High German \emph{zorn} (EtymDB), already Old High German \emph{zorn}; its older origin is not settled.
\end{cousinweb}

\begin{grammarlens}[title={What you've built}]
One more word for telling, explaining and describing at B1.
\end{grammarlens}

\subsection*{Your turn}
\begin{itemize}
\raggedright
  \item \emph{Say it:} \emph{der Zorn}
  \item \emph{Say it:} \emph{der Zorn}, once more
  \item \emph{Say it:} a sentence of your own with \emph{der Zorn}
  \item \emph{Recall:} say the German for to blush, then the German for rage, then say \emph{der Zorn} again
\end{itemize}

\begin{tcolorbox}[breakable,colback=teal!4,colframe=teal!35!black,title={Before you move on}]
What does der Zorn mean? (\textquotedblleft{}Anger, wrath\textquotedblright{}.) Say it once more.
\end{tcolorbox}
\section[ärgern]{ärgern — to annoy; to be annoyed (sich ärgern)}

\label{lesson:GE-C468-aergern}



\begin{quote}
Before the new one: say the German for rage, then the German for anger.
\end{quote}

\subsection*{You'll want to know: ärgern}
\textbf{ärgern} — \textquotedblleft{}to annoy; to be annoyed (sich ärgern)\textquotedblright{}.

\begin{quote}
Ich ärgere mich über den Lärm der Nachbarn.
\end{quote}

\begin{quote}
\emph{I'm annoyed about the neighbours' noise.}
\end{quote}

\emph{Sich ärgern über} + accusative — to be annoyed about. Without \emph{sich}: \emph{Du ärgerst mich!} — You're annoying me!

\begin{cousinweb}
\emph{ärgern} is inherited from Old High German \emph{argerōn} \textquotedblleft{}to make worse\textquotedblright{} (EtymDB), built on \emph{arg} \textquotedblleft{}bad\textquotedblright{}, from Proto-Germanic \emph{argaz}.
\end{cousinweb}

\begin{grammarlens}[title={What you've built}]
One more word for telling, explaining and describing at B1.
\end{grammarlens}

\subsection*{Your turn}
\begin{itemize}
\raggedright
  \item \emph{Say it:} \emph{ärgern}
  \item \emph{Say it:} \emph{ärgern}, once more
  \item \emph{Say it:} a sentence of your own with \emph{ärgern}
  \item \emph{Recall:} say the German for rage, then the German for anger, then say \emph{ärgern} again
\end{itemize}

\begin{tcolorbox}[breakable,colback=teal!4,colframe=teal!35!black,title={Before you move on}]
What does ärgern mean? (\textquotedblleft{}To annoy; to be annoyed (sich ärgern)\textquotedblright{}.) Say it once more.
\end{tcolorbox}
\section[aufregen]{aufregen — to get worked up (sich aufregen); to upset}

\label{lesson:GE-C468-aufregen}



\begin{quote}
Before the new one: say the German for anger, then the German for to annoy.
\end{quote}

\subsection*{You'll want to know: aufregen}
\textbf{aufregen} — \textquotedblleft{}to get worked up (sich aufregen); to upset\textquotedblright{}.

\begin{quote}
Mein Opa regt sich über jede Kleinigkeit auf.
\end{quote}

\begin{quote}
\emph{My grandpa gets worked up about every little thing.}
\end{quote}

Separable and mostly reflexive: \emph{sich über etwas aufregen}. \emph{Das regt mich auf!} — That makes me mad!

\begin{cousinweb}
\emph{aufregen} is built on \emph{auf} \textquotedblleft{}up\textquotedblright{} and \emph{regen} \textquotedblleft{}to stir, to move\textquotedblright{}, which goes back to Middle High German \emph{regen} \textquotedblleft{}to excite, to move\textquotedblright{} (EtymDB).
\end{cousinweb}

\begin{grammarlens}[title={What you've built}]
One more word for telling, explaining and describing at B1.
\end{grammarlens}

\subsection*{Your turn}
\begin{itemize}
\raggedright
  \item \emph{Say it:} \emph{aufregen}
  \item \emph{Say it:} \emph{aufregen}, once more
  \item \emph{Say it:} a sentence of your own with \emph{aufregen}
  \item \emph{Recall:} say the German for anger, then the German for to annoy, then say \emph{aufregen} again
\end{itemize}

\begin{tcolorbox}[breakable,colback=teal!4,colframe=teal!35!black,title={Before you move on}]
What does aufregen mean? (\textquotedblleft{}To get worked up (sich aufregen); to upset\textquotedblright{}.) Say it once more.
\end{tcolorbox}
\section[toben]{toben — to rage; to romp}

\label{lesson:GE-C468-toben}



\begin{quote}
Before the new one: say the German for to annoy, then the German for to get worked up.
\end{quote}

\subsection*{You'll want to know: toben}
\textbf{toben} — \textquotedblleft{}to rage; to romp\textquotedblright{}.

\begin{quote}
Die Kinder toben im Garten.
\end{quote}

\begin{quote}
\emph{The children are romping in the garden.}
\end{quote}

Of people in anger, of storms, and of children playing wildly.

\begin{cousinweb}
The verb \emph{toben} goes back to Old High German \emph{tobēn} \textquotedblleft{}to rage, to be mad\textquotedblright{}; its older origin is not settled, and EtymDB has no entry for it.
\end{cousinweb}

\begin{grammarlens}[title={What you've built}]
One more word for telling, explaining and describing at B1.
\end{grammarlens}

\subsection*{Your turn}
\begin{itemize}
\raggedright
  \item \emph{Say it:} \emph{toben}
  \item \emph{Say it:} \emph{toben}, once more
  \item \emph{Say it:} a sentence of your own with \emph{toben}
  \item \emph{Recall:} say the German for to annoy, then the German for to get worked up, then say \emph{toben} again
\end{itemize}

\begin{tcolorbox}[breakable,colback=teal!4,colframe=teal!35!black,title={Before you move on}]
What does toben mean? (\textquotedblleft{}To rage; to romp\textquotedblright{}.) Say it once more.
\end{tcolorbox}
